\paragraph{隐藏位压力微分、Bernoulli 递阶反演与谱零块的二项式因子化}
在推前熵账本、Rényi 散度冻结、规范群体积的 Stirling--Bernoulli 层级，以及约数驱动的谱零块均已建立之后，还可把隐藏位统计、碰撞指纹、Bernoulli--\(\zeta\) 反演与群环多项式的代数刚性进一步压缩为下列一组闭合结论。

\begin{theorem}[压力谱在 \texorpdfstring{$a=1$}{a=1} 的一、二阶导数分别给出隐藏对数多重度的密度与涨落]\label{thm:xi-time-part60ab-hidden-logmultiplicity-pressure-derivatives}
\leanverified{paper\_xi\_time\_part60ab\_hidden\_logmultiplicity\_pressure\_derivatives}
沿用 H.3.1 的纤维多重度记号。对每个 \(m\) 记
\[
\mu_m(x):=\frac{d_m(x)}{2^m},
\qquad
S_a(m):=\sum_{x\in X_m} d_m(x)^a,
\qquad
P_m(a):=\frac1m\log S_a(m).
\]
设开区间 \(I\subset (0,\infty)\) 满足 \(1\in I\)，并假设每个 \(P_m\) 在 \(I\) 上属于 \(C^2\)，且存在 \(P\in C^2(I)\) 使
\[
P_m\longrightarrow P
\qquad\text{于 }C^2_{\mathrm{loc}}(I).
\]
若 \(X\sim \mu_m\)，定义隐藏对数多重度随机变量
\[
\mathscr H_m:=\log d_m(X).
\]
则对每个 \(m\) 都有精确恒等式
\[
\frac1m\EE_{\mu_m}[\mathscr H_m]=P_m'(1),
\qquad
\frac1m\Var_{\mu_m}(\mathscr H_m)=P_m''(1),
\]
从而
\[
\boxed{\;
\lim_{m\to\infty}\frac1m\EE_{\mu_m}[\mathscr H_m]=P'(1),\;}
\]
\[
\boxed{\;
\lim_{m\to\infty}\frac1m\Var_{\mu_m}(\mathscr H_m)=P''(1).\;}
\]
\end{theorem}
\begin{proof}
定义
\[
\Phi_m(a):=\frac1m\log\sum_{x\in X_m}\mu_m(x)^a.
\]
由
\[
\mu_m(x)=\frac{d_m(x)}{2^m}
\]
可得
\[
\Phi_m(a)
=
\frac1m\log\sum_x d_m(x)^a-a\log 2
=
P_m(a)-a\log 2.
\]

另一方面，
\[
\Phi_m'(a)
=
\frac1m
\frac{\sum_x \mu_m(x)^a\log \mu_m(x)}{\sum_x\mu_m(x)^a}.
\]
在 \(a=1\) 处利用 \(\sum_x\mu_m(x)=1\)，得到
\[
\Phi_m'(1)
=
\frac1m\sum_x\mu_m(x)\log\mu_m(x)
=
\frac1m\sum_x\mu_m(x)\log d_m(x)-\log 2
=
\frac1m\EE_{\mu_m}[\mathscr H_m]-\log 2.
\]
而由 \(\Phi_m(a)=P_m(a)-a\log 2\) 又有
\[
\Phi_m'(1)=P_m'(1)-\log 2.
\]
两式相减即得
\[
\frac1m\EE_{\mu_m}[\mathscr H_m]=P_m'(1).
\]

再看二阶导数。标准指数族计算给出
\[
\Phi_m''(1)
=
\frac1m
\left(
\sum_x\mu_m(x)(\log\mu_m(x))^2
\;-\;
\Bigl(\sum_x\mu_m(x)\log\mu_m(x)\Bigr)^2
\right)
=
\frac1m\Var_{\mu_m}(\log\mu_m(X)).
\]
又因
\[
\log\mu_m(X)=\mathscr H_m-m\log 2,
\]
加常数不改变方差，故
\[
\Var_{\mu_m}(\log\mu_m(X))=\Var_{\mu_m}(\mathscr H_m).
\]
另一方面，
\[
\Phi_m''(1)=P_m''(1).
\]
遂得
\[
\frac1m\Var_{\mu_m}(\mathscr H_m)=P_m''(1).
\]

最后，由 \(C^2_{\mathrm{loc}}(I)\) 收敛，必有
\[
P_m'(1)\to P'(1),
\qquad
P_m''(1)\to P''(1).
\]
代入前述有限尺度恒等式即得极限公式。证毕。
\end{proof}

\begin{theorem}[隐藏位缺口的 KL 识别与 Rényi--\texorpdfstring{$2$}{2} 碰撞证书]\label{thm:xi-time-part60ab-hidden-gap-kl-renyi2-collision}
沿用表 \ref{tab:fold_multiplicity_histogram} 的隐藏位预算记号。记
\[
u_m(x):=\frac1{|X_m|}
\qquad (x\in X_m),
\]
并定义隐藏位、可见位与容量缺口
\[
h_m:=\EE_{\mu_m}\bigl[\log_2 d_m(X)\bigr],
\qquad
I_m:=m-h_m,
\qquad
\mathrm{gap}_m:=\log_2|X_m|-I_m.
\]
再定义二阶碰撞指纹
\[
C_{\mathrm{col}}(m):=
|X_m|\sum_{x\in X_m}\mu_m(x)^2.
\]
则有严格恒等式
\[
\boxed{\;
\mathrm{gap}_m
=
\frac{1}{\log 2}\,
D_{\mathrm{KL}}(\mu_m\|u_m),\;}
\]
以及精确上界
\[
\boxed{\;
\mathrm{gap}_m
\le
\log_2 C_{\mathrm{col}}(m).\;}
\]
并且下列三条等价：
\begin{enumerate}
  \item \(\mathrm{gap}_m=\log_2 C_{\mathrm{col}}(m)\)；
  \item \(\mu_m=u_m\)；
  \item \(d_m(x)\) 在 \(X_m\) 上常值。
\end{enumerate}
\end{theorem}
\begin{proof}
由定义，
\[
\mathrm{gap}_m
=
\log_2|X_m|-m+\sum_{x\in X_m}\mu_m(x)\log_2 d_m(x).
\]
再用
\[
\mu_m(x)=\frac{d_m(x)}{2^m},
\qquad
u_m(x)=\frac1{|X_m|},
\]
得到
\[
\mathrm{gap}_m
=
\sum_x \mu_m(x)\log_2\frac{d_m(x)|X_m|}{2^m}
=
\sum_x \mu_m(x)\log_2\frac{\mu_m(x)}{u_m(x)}
=
\frac{1}{\log 2}\,D_{\mathrm{KL}}(\mu_m\|u_m).
\]
这便给出第一条恒等式。

另一方面，Rényi 散度对阶数单调，故
\[
D_{\mathrm{KL}}(\mu_m\|u_m)\le D_2(\mu_m\|u_m).
\]
而
\[
D_2(\mu_m\|u_m)
=
\log\sum_x \frac{\mu_m(x)^2}{u_m(x)}
=
\log\!\Bigl(|X_m|\sum_x \mu_m(x)^2\Bigr)
=
\log C_{\mathrm{col}}(m).
\]
两边除以 \(\log 2\) 即得
\[
\mathrm{gap}_m\le \log_2 C_{\mathrm{col}}(m).
\]

最后，\(D_{\mathrm{KL}}(\mu_m\|u_m)=D_2(\mu_m\|u_m)\) 当且仅当对 \(\mu_m\)-几乎处处的 \(x\)，比值 \(\mu_m(x)/u_m(x)\) 为常数。由于 \(u_m\) 在 \(X_m\) 上严格正、且二者都是概率分布，这等价于 \(\mu_m=u_m\)。而
\[
\mu_m=u_m
\iff
\frac{d_m(x)}{2^m}=\frac1{|X_m|}
\ \ (\forall x\in X_m)
\iff
d_m(x)\ \text{在 }X_m\text{ 上常值}.
\]
证毕。
\end{proof}

\begin{theorem}[fold-gauge 常数列的逐阶极限消元反演给出全部偶 Bernoulli 数与偶值 \texorpdfstring{$\zeta(2r)$}{zeta(2r)}]\label{thm:xi-time-part60ab-gauge-constant-recursive-bernoulli-stripping}
沿用定理 \ref{thm:fold-bin-gauge-constant-stirling-bernoulli-hierarchy} 的记号，记
\[
q:=\frac{\varphi}{2}\in(0,1),
\qquad
E_m:=\log C_m-\log(2\pi).
\]
并对每个整数 \(r\ge 1\) 定义
\[
a_r:=
\frac{B_{2r}}{r(2r-1)}
\bigl(\varphi^{-1}+\varphi^{2r-3}\bigr).
\]
再递归定义残差
\[
R_m^{(1)}:=E_m,
\qquad
R_m^{(r)}:=R_m^{(r-1)}-a_{r-1}q^{(2r-3)m}
\qquad (r\ge 2).
\]
则对每个 \(r\ge 1\) 都有极限
\[
\boxed{\;
\lim_{m\to\infty}q^{-(2r-1)m}R_m^{(r)}=a_r.\;}
\]
从而
\[
\boxed{\;
B_{2r}
=
\frac{r(2r-1)}{\varphi^{-1}+\varphi^{2r-3}}
\lim_{m\to\infty}
\Bigl(\frac{2}{\varphi}\Bigr)^{(2r-1)m}R_m^{(r)}.\;}
\]
\[
\boxed{\;
\zeta(2r)
=
(-1)^{r-1}\frac{(2\pi)^{2r}}{2(2r)!}\,
\frac{r(2r-1)}{\varphi^{-1}+\varphi^{2r-3}}
\lim_{m\to\infty}
\Bigl(\frac{2}{\varphi}\Bigr)^{(2r-1)m}R_m^{(r)}.\;}
\]
因此，单一标量序列 \((\log C_m)_{m\ge 1}\) 已足以逐阶恢复全部偶 Bernoulli 数与全部偶值 \(\zeta(2r)\)。
\end{theorem}
\begin{proof}
由定理 \ref{thm:fold-bin-gauge-constant-stirling-bernoulli-hierarchy}，对任意固定 \(R\ge 1\) 有
\[
E_m
=
\sum_{s=1}^{R}a_s q^{(2s-1)m}
+O\!\bigl(q^{(2R+1)m}\bigr).
\]
对上式取 \(R=r\)，并逐次减去前 \(r-1\) 个主项，得到
\[
R_m^{(r)}
=
a_r q^{(2r-1)m}+O\!\bigl(q^{(2r+1)m}\bigr).
\]
两边同乘 \(q^{-(2r-1)m}\) 得
\[
q^{-(2r-1)m}
R_m^{(r)}
=
a_r+O(q^{2m}),
\]
令 \(m\to\infty\) 即得极限公式。
再由
\[
a_r=
\frac{B_{2r}}{r(2r-1)}
\bigl(\varphi^{-1}+\varphi^{2r-3}\bigr)
\]
即可反解出 \(B_{2r}\)。
再由标准 Bernoulli--\(\zeta\) 恒等式
\[
B_{2r}
=
(-1)^{r-1}\frac{2(2r)!}{(2\pi)^{2r}}\zeta(2r),
\]
可立即反解出 \(\zeta(2r)\) 的显式恢复公式。证毕。
\end{proof}

\begin{corollary}[fold-gauge 常数的唯一主模与比例极限]\label{cor:xi-time-part60ab-gauge-constant-leading-mode-ratio}
沿用定理 \ref{thm:xi-time-part60ab-gauge-constant-recursive-bernoulli-stripping} 的记号，则
\[
\boxed{\;
\log C_m-\log(2\pi)
=
a_1 q^m+O(q^{3m}),
\qquad
q=\frac{\varphi}{2}.\;}
\]
从而
\[
\boxed{\;
\lim_{m\to\infty}
\frac{\log C_{m+1}-\log(2\pi)}
{\log C_m-\log(2\pi)}
=
\frac{\varphi}{2},\;}
\]
以及
\[
\boxed{\;
\lim_{m\to\infty}
\Bigl(\frac{2}{\varphi}\Bigr)^m
\bigl(\log C_m-\log(2\pi)\bigr)
=
a_1.\;}
\]
\end{corollary}
\begin{proof}
由上一条定理在 \(r=1\) 时的系数公式，
\[
a_1=
\frac{B_2}{1\cdot 1}
\bigl(\varphi^{-1}+\varphi^{-1}\bigr)
=
\frac16\cdot 2\varphi^{-1}
=
\frac{1}{3\varphi}.
\]
故
\[
\log C_m-\log(2\pi)
=
a_1 q^m+O(q^{3m}).
\]
于是
\[
\frac{\log C_{m+1}-\log(2\pi)}
{\log C_m-\log(2\pi)}
=
\frac{\frac{1}{3\varphi}q^{m+1}+O(q^{3m+3})}
{\frac{1}{3\varphi}q^m+O(q^{3m})}
\longrightarrow q=\frac{\varphi}{2},
\]
并且
\[
\Bigl(\frac{2}{\varphi}\Bigr)^m\bigl(\log C_m-\log(2\pi)\bigr)
=
\frac{1}{3\varphi}+O(q^{2m})
\longrightarrow
\frac{1}{3\varphi}.
\]
证毕。
\end{proof}

在上述逐阶极限消元公式之后，\(\log C_m\) 的 Bernoulli--\(\zeta\) 谱链还可进一步压缩为以下三条结构后果。

\begin{theorem}[fold-gauge 常数的 Richardson--Neville 任意阶消模器与 \texorpdfstring{$(2R+1)$}{(2R+1)} 级加速]\label{thm:xi-time-part60ab-gauge-constant-richardson-neville-accelerator}
沿用定理 \ref{thm:xi-time-part60ab-gauge-constant-recursive-bernoulli-stripping} 的记号，并定义
\[
\lambda_r:=q^{2r-1}
\qquad (r\ge 1).
\]
对任意固定整数 \(R\ge 1\)，令
\[
Q_R(z):=\prod_{r=1}^{R}\frac{z-\lambda_r}{1-\lambda_r}
=
\sum_{j=0}^{R}\gamma_{R,j}z^j,
\]
并定义线性加速器
\[
\mathcal A_R(m):=\sum_{j=0}^{R}\gamma_{R,j}\log C_{m+j}.
\]
则有
\[
\boxed{\;
\mathcal A_R(m)=\log(2\pi)+O\!\bigl(q^{(2R+1)m}\bigr).\;}
\]
换言之，\(\mathcal A_R\) 精确消去 \(\log C_m\) 渐近展开中的前 \(R\) 个 Bernoulli 模态，并把逼近 \(\log(2\pi)\) 的误差阶提升到 \(q^{(2R+1)m}\)。
\end{theorem}
\begin{proof}
由定义立即有
\[
Q_R(1)=1,
\qquad
Q_R(\lambda_r)=0
\qquad (1\le r\le R).
\]
另一方面，由定理 \ref{thm:fold-bin-gauge-constant-stirling-bernoulli-hierarchy}，对固定 \(R\) 有
\[
\log C_{m+j}
=
\log(2\pi)+\sum_{r=1}^{R}a_r\lambda_r^{m+j}
+O\!\bigl(\lambda_{R+1}^{m+j}\bigr)
\qquad (0\le j\le R).
\]
由于 \(j\) 仅在有限集合 \(\{0,\dots,R\}\) 中变化，上式余项可统一写为 \(O(\lambda_{R+1}^m)\)。将其代入 \(\mathcal A_R(m)\)，得到
\[
\mathcal A_R(m)
=
\log(2\pi)\sum_{j=0}^{R}\gamma_{R,j}
+\sum_{r=1}^{R}a_r\lambda_r^m\sum_{j=0}^{R}\gamma_{R,j}\lambda_r^j
+O\!\bigl(\lambda_{R+1}^m\bigr).
\]
再由
\[
\sum_{j=0}^{R}\gamma_{R,j}=Q_R(1)=1,
\qquad
\sum_{j=0}^{R}\gamma_{R,j}\lambda_r^j=Q_R(\lambda_r)=0
\qquad (1\le r\le R),
\]
主项中的前 \(R\) 个 Bernoulli 模态全部被精确消去，从而
\[
\mathcal A_R(m)
=
\log(2\pi)+O\!\bigl(\lambda_{R+1}^m\bigr)
=
\log(2\pi)+O\!\bigl(q^{(2R+1)m}\bigr).
\]
证毕。
\end{proof}

\begin{theorem}[fold-gauge 误差序列的非 \texorpdfstring{$C$}{C}-有限性与普通生成函数的非有理性]\label{thm:xi-time-part60ab-gauge-constant-non-cfinite-ogf-obstruction}
\leanverified{paper\_xi\_time\_part60ab\_gauge\_constant\_non\_cfinite\_ogf\_obstruction}
令
\[
u_m:=\log C_m-\log(2\pi).
\]
则序列 \((u_m)_{m\ge 1}\) 不是任何常系数线性递推序列，即不是 \(C\)-有限序列。因而其普通生成函数
\[
U(z):=\sum_{m\ge 1}u_m z^m
\]
不可能是有理函数。
\end{theorem}
\begin{proof}
反设 \((u_m)_{m\ge 1}\) 为 \(C\)-有限序列。则存在有限个复数 \(\mu_1,\dots,\mu_s\) 与多项式 \(P_1,\dots,P_s\)，使得对充分大的 \(m\) 有
\[
u_m=\sum_{j=1}^{s}P_j(m)\mu_j^m.
\]
因此 \(u_m\) 只可能由有限多个指数底数 \(\mu_j\) 生成。

另一方面，由定理 \ref{thm:xi-time-part60ab-gauge-constant-recursive-bernoulli-stripping}，对每个 \(R\ge 1\) 都有
\[
u_m-\sum_{r=1}^{R-1}a_r q^{(2r-1)m}
=
a_R q^{(2R-1)m}+O\!\bigl(q^{(2R+1)m}\bigr).
\]
又由于
\[
a_R=
\frac{B_{2R}}{R(2R-1)}
\bigl(\varphi^{-1}+\varphi^{2R-3}\bigr)
\]
且 \(B_{2R}\neq 0\)，故 \(a_R\neq 0\) 对每个 \(R\ge 1\) 都成立。于是对每个 \(R\)，尺度 \(q^{(2R-1)m}\) 都真实出现在 \(u_m\) 的渐近链中。另一方面，
\[
u_m-\sum_{r=1}^{R-1}a_r q^{(2r-1)m}
\]
仍是由有限个指数底数
\[
\{\mu_1,\dots,\mu_s,q,q^3,\dots,q^{2R-3}\}
\]
生成的多项式--指数和。由其渐近主项为非零的 \(a_R q^{(2R-1)m}\)，可知底数 \(q^{2R-1}\) 必属于上式有限集合。由于
\[
q^{2R-1}\notin \{q,q^3,\dots,q^{2R-3}\},
\]
遂只能有
\[
q^{2R-1}\in \{\mu_1,\dots,\mu_s\}.
\]
这与 \(\{\mu_1,\dots,\mu_s\}\) 有限而
\[
q,q^3,q^5,\dots
\]
两两不同矛盾。故 \((u_m)\) 不是 \(C\)-有限序列。

最后，普通生成函数有理当且仅当其系数序列 \(C\)-有限，这是经典等价命题，故 \(U(z)\) 不可能为有理函数。证毕。
\end{proof}

\begin{corollary}[fold-gauge Bernoulli 塔的无限模态障碍]\label{cor:xi-time-part60ab-gauge-constant-infinite-modal-obstruction}
令
\[
E_m:=\log C_m-\log(2\pi).
\]
则不存在有限组互异复数 \(\lambda_1,\dots,\lambda_s\) 与多项式 \(P_1,\dots,P_s\)，使得对充分大的 \(m\) 恒有
\[
E_m=\sum_{j=1}^{s}P_j(m)\lambda_j^m.
\]
等价地，H.111 的 Stirling--Bernoulli 层级展开对应一条真正的无限模态谱链，而非任何有限维 companion 机制所能封装的有限指数叠加。
\end{corollary}
\begin{proof}
这正是定理 \ref{thm:xi-time-part60ab-gauge-constant-non-cfinite-ogf-obstruction} 的直接重述，因为常系数线性递推序列与有限个多项式--指数项的和完全等价。证毕。
\end{proof}

\begin{theorem}[高阶尾部精度强制前导 Bernoulli--\texorpdfstring{$\zeta$}{zeta} 模态一致]\label{thm:xi-time-part60ab-gauge-constant-tail-rigidity}
沿用定理 \ref{thm:xi-time-part60ab-gauge-constant-recursive-bernoulli-stripping} 的记号。设另一实序列 \((D_m)_{m\ge 1}\) 满足：对某个固定整数 \(R\ge 1\)，存在实系数 \(b_1,\dots,b_R\) 使
\[
D_m-\log(2\pi)
=
\sum_{r=1}^{R}b_r q^{(2r-1)m}
+O\!\bigl(q^{(2R+1)m}\bigr),
\]
并且同时有
\[
D_m-\log C_m
=
O\!\bigl(q^{(2R+1)m}\bigr).
\]
则必有
\[
b_r=a_r
\qquad (1\le r\le R).
\]
因而前 \(R\) 个偶 Bernoulli 数 \(B_{2r}\) 与偶值 \(\zeta(2r)\) 也被同一尾部精度唯一锁定。
\end{theorem}
\begin{proof}
由定理 \ref{thm:fold-bin-gauge-constant-stirling-bernoulli-hierarchy}，
\[
\log C_m-\log(2\pi)
=
\sum_{r=1}^{R}a_r q^{(2r-1)m}
+O\!\bigl(q^{(2R+1)m}\bigr).
\]
将其与 \(D_m\) 的展开相减，得到
\[
D_m-\log C_m
=
\sum_{r=1}^{R}(b_r-a_r)q^{(2r-1)m}
+O\!\bigl(q^{(2R+1)m}\bigr).
\]
若存在最小的 \(r_0\le R\) 使 \(b_{r_0}\neq a_{r_0}\)，则上式可改写为
\[
D_m-\log C_m
=
(b_{r_0}-a_{r_0})q^{(2r_0-1)m}
+O\!\bigl(q^{(2r_0+1)m}\bigr).
\]
两边同乘 \(q^{-(2r_0-1)m}\)，并令 \(m\to\infty\)，便得到
\[
q^{-(2r_0-1)m}\bigl(D_m-\log C_m\bigr)\to b_{r_0}-a_{r_0}\neq 0,
\]
这与假设
\[
D_m-\log C_m
=
O\!\bigl(q^{(2R+1)m}\bigr)
\]
矛盾，因为后者特别蕴含
\(
D_m-\log C_m=O(q^{(2r_0+1)m})
\)。
故 \(b_r=a_r\) 对全部 \(1\le r\le R\) 成立。

最后，由定理 \ref{thm:xi-time-part60ab-gauge-constant-recursive-bernoulli-stripping} 中 \(a_r\) 与 \(B_{2r}\)、\(\zeta(2r)\) 的显式对应关系，前 \(R\) 个 Bernoulli--\(\zeta\) 模态也随之被唯一锁定。证毕。
\end{proof}

\endinput


\paragraph{Stirling--Bernoulli 模态的移位投影与偶值 \texorpdfstring{$\zeta$}{zeta} 显式恢复}
附录中的 Bernoulli 层级、递归抽取与有限阶移位消元已经给出三条彼此兼容的读出接口。把它们压回正文主线，可得到如下三条直接后果：每个 Stirling--Bernoulli 模态都可由一个显式有限阶移位投影子直接抽出；偶值 \(\zeta(2r)\) 因而成为单一规范常数序列上的直接可审计读出；而无限多个互异衰减模的同时存在，则排除了任何有限 companion 型封装。

\begin{theorem}[fold-gauge Bernoulli 模态的有限阶移位投影公式]\label{thm:xi-time-part60ab2-logcm-shift-projector}
\leanverified{paper\_xi\_time\_part60ab2\_logcm\_shift\_projector}
沿用定理 \ref{thm:fold-bin-gauge-constant-stirling-bernoulli-hierarchy} 的记号，记
\[
f_m:=\log C_m-\log(2\pi),
\qquad
q_r:=\left(\frac{\varphi}{2}\right)^{2r-1},
\qquad
a_r:=\frac{B_{2r}}{r(2r-1)}\bigl(\varphi^{-1}+\varphi^{2r-3}\bigr).
\]
并令移位算子 \(E\) 作用于序列 \(f=(f_m)_{m\ge 1}\) 为
\[
(Ef)_m:=f_{m+1}.
\]
对每个 \(r\ge 1\) 定义
\[
\mathcal P_{r-1}(E):=\prod_{j=1}^{r-1}(E-q_j),
\]
其中 \(r=1\) 时取空积 \(\mathcal P_0(E)=1\)。则极限
\[
\lim_{m\to\infty}
\frac{q_r^{-m}\bigl(\mathcal P_{r-1}(E)f\bigr)_m}
{\prod_{j=1}^{r-1}(q_r-q_j)}
\]
存在，并且恰等于 \(a_r\)。特别地，每个模态系数 \(a_r\) 都可由 \((\log C_m)_{m\ge 1}\) 的一个有限阶移位投影极限直接读出。
\end{theorem}
\begin{proof}
定理 \ref{thm:derived-fold-gauge-shift-annihilator-even-zeta} 已给出
\[
\lim_{m\to\infty}q_r^{-m}\bigl(\mathcal P_{r-1}(E)f\bigr)_m
=
a_r\prod_{j=1}^{r-1}(q_r-q_j).
\]
又由于
\[
q_1>q_2>\cdots>0,
\]
故对每个 \(j<r\) 都有 \(q_r-q_j\neq 0\)。两边同除以该非零乘积，即得所示公式。证毕。
\end{proof}

\begin{corollary}[有限阶移位投影给出的偶值 \texorpdfstring{$\zeta(2r)$}{zeta(2r)} 显式恢复公式]\label{cor:xi-time-part60ab2-logcm-shift-even-zeta}
\leanverified{paper\_xi\_time\_part60ab2\_logcm\_shift\_even\_zeta}
沿用定理 \ref{thm:xi-time-part60ab2-logcm-shift-projector} 的记号。对每个 \(r\ge 1\)，都有
\[
\zeta(2r)
=
(-1)^{r-1}\frac{(2\pi)^{2r}}{2(2r)!}\,
\frac{r(2r-1)}{\varphi^{-1}+\varphi^{2r-3}}\,
\frac{
\lim\limits_{m\to\infty}
q_r^{-m}\bigl(\mathcal P_{r-1}(E)f\bigr)_m
}
{\prod_{j=1}^{r-1}(q_r-q_j)}.
\]
因此，偶值 \(\zeta\)-谱并不只作为 \((\log C_m)\) 的渐近尾项存在，而是可由一族显式有限阶移位投影子逐模态直接恢复。
\end{corollary}
\begin{proof}
由定理 \ref{thm:xi-time-part60ab2-logcm-shift-projector}，
\[
a_r=
\frac{
\lim\limits_{m\to\infty}
q_r^{-m}\bigl(\mathcal P_{r-1}(E)f\bigr)_m
}
{\prod_{j=1}^{r-1}(q_r-q_j)}.
\]
再代入
\[
a_r=\frac{B_{2r}}{r(2r-1)}\bigl(\varphi^{-1}+\varphi^{2r-3}\bigr),
\qquad
B_{2r}=(-1)^{r-1}\frac{2(2r)!}{(2\pi)^{2r}}\zeta(2r),
\]
整理即得。证毕。
\end{proof}

\begin{corollary}[归一化移位投影子的单步 Bernoulli 恢复]\label{cor:xi-time-part60ab2-logcm-normalized-projector-bernoulli}
\leanverified{paper\_xi\_time\_part60ab2\_logcm\_normalized\_projector\_bernoulli}
沿用定理 \ref{thm:xi-time-part60ab2-logcm-shift-projector} 的记号。对每个 \(r\ge 1\)，定义归一化移位投影子
\[
\Pi_r(E):=
\prod_{j=1}^{r-1}\frac{E-q_j}{q_r-q_j},
\]
其中 \(r=1\) 时取空积 \(\Pi_1(E)=1\)。
则
\[
\lim_{m\to\infty}q_r^{-m}\bigl(\Pi_r(E)f\bigr)_m=a_r,
\]
从而
\[
\lim_{m\to\infty}q_r^{-m}\bigl(\Pi_r(E)f\bigr)_m
=
\frac{B_{2r}}{r(2r-1)}\bigl(\varphi^{-1}+\varphi^{2r-3}\bigr),
\]
以及
\[
B_{2r}
=
\frac{r(2r-1)}{\varphi^{-1}+\varphi^{2r-3}}
\lim_{m\to\infty}q_r^{-m}\bigl(\Pi_r(E)f\bigr)_m.
\]
因此，偶 Bernoulli 数可由一族封闭形式的归一化移位投影子直接逐阶读出，而无需求助于递归剥离此前各级模态。
\end{corollary}
\begin{proof}
由定义，
\[
\Pi_r(E)
=
\frac{\mathcal P_{r-1}(E)}{\prod_{j=1}^{r-1}(q_r-q_j)}.
\]
故定理 \ref{thm:xi-time-part60ab2-logcm-shift-projector} 立即给出
\[
\lim_{m\to\infty}q_r^{-m}\bigl(\Pi_r(E)f\bigr)_m=a_r.
\]
再代入
\[
a_r=\frac{B_{2r}}{r(2r-1)}\bigl(\varphi^{-1}+\varphi^{2r-3}\bigr)
\]
并整理，即得所示 Bernoulli 恢复公式。证毕。
\end{proof}

\begin{theorem}[有限阶移位投影强制 fold-gauge 谱塔无限]\label{thm:xi-time-part60ab2-logcm-shift-infinite-modal}
序列
\[
f_m=\log C_m-\log(2\pi)
\]
不是任何有限阶常系数线性递推序列；等价地，其普通生成函数
\[
F(z):=\sum_{m\ge 1}f_m z^m
\]
不是有理函数。
\end{theorem}
\leanverified{paper\_xi\_time\_part60ab2\_logcm\_shift\_infinite\_modal}
\begin{proof}
由定理 \ref{thm:xi-time-part60ab2-logcm-shift-projector}，对每个 \(r\ge 1\) 都有
\[
\lim_{m\to\infty}
\frac{q_r^{-m}\bigl(\mathcal P_{r-1}(E)f\bigr)_m}
{\prod_{j=1}^{r-1}(q_r-q_j)}
=a_r.
\]
而
\[
a_r=
\frac{B_{2r}}{r(2r-1)}\bigl(\varphi^{-1}+\varphi^{2r-3}\bigr)\neq 0,
\]
因为偶 Bernoulli 数 \(B_{2r}\) 从不消失，且 \(\varphi^{-1}+\varphi^{2r-3}>0\)。因此，\(f_m\) 的渐近谱包含无限多个互异底数
\[
q_r=\left(\frac{\varphi}{2}\right)^{2r-1},
\qquad
r=1,2,\dots.
\]
另一方面，任何有限阶常系数线性递推序列的渐近谱都只能由有限多个指数底数构成，故不可能容纳上述无限模态塔。于是 \(f_m\) 非 \(C\)-有限。

普通生成函数有理当且仅当其系数序列 \(C\)-有限，这是经典等价命题，遂得 \(F(z)\) 非有理。证毕。
\end{proof}

\noindent
上述三条结果把 \((\log C_m)\) 的 Bernoulli 层级从“逐阶剥离尾项”的恢复程序，提升为一族显式有限阶移位滤波器的直接读出接口。因而，偶值 \(\zeta\)-谱并非只是在渐近展开中被动出现，而可由单一规范常数序列经过确定性的有限阶投影逐模态抽出；同时，无限模态障碍又说明这种读出结构本质上超出任何有限状态线性递推的封装能力。

\endinput


\begin{theorem}[约数触发的谱零块强制群环多项式含有显式二项式因子]\label{thm:xi-time-part60ab-zero-block-forces-binomial-factor}
\leanverified{paper\_xi\_time\_part60ab\_zero\_block\_forces\_binomial\_factor}
沿用命题 \ref{prop:fold-multiplicity-group-algebra} 的记号。记
\[
M:=F_{m+2},
\qquad
D_m(x):=\sum_{r=0}^{M-1}d_m(r)x^r\in\ZZ[x],
\qquad
\deg D_m<M.
\]
若 \(d\mid (m+2)\) 是真约数且满足
\[
F_d\mid \frac{M}{2},
\]
则有整系数因子结论
\[
\boxed{\;
x^{F_d}+1\mid D_m(x)\qquad\text{于 }\ZZ[x].\;}
\]
特别地，当
\[
m+2\equiv 0\pmod 6
\]
并取 \(d=(m+2)/2\) 时，
\[
\boxed{\;
x^{F_{(m+2)/2}}+1\mid D_m(x).\;}
\]
\end{theorem}
\begin{proof}
令
\[
t:=\frac{M}{2F_d}\in\NN.
\]
由推论 \ref{cor:fold-zero-divisor-lb}，
\[
S_{F_d}(M)\subseteq \mathcal Z_m
:=
\{k\in\ZZ/(M\ZZ):\ \widehat d_m(k)=0\}.
\]
对每个 \(\ell=0,1,\dots,F_d-1\)，取
\[
k_\ell:=(2\ell+1)t.
\]
则
\[
k_\ell F_d=(2\ell+1)\frac{M}{2}\equiv \frac{M}{2}\pmod M,
\]
故 \(k_\ell\in S_{F_d}(M)\subseteq \mathcal Z_m\)。于是由离散 Fourier 变换的定义，
\[
D_m(\omega_M^{-k_\ell})=\widehat d_m(k_\ell)=0,
\qquad
\omega_M:=e^{2\pi i/M}.
\]

另一方面，
\[
(\omega_M^{-k_\ell})^{F_d}
=
\exp\!\left(-2\pi i\,\frac{k_\ell F_d}{M}\right)
=
\exp\!\left(-(2\ell+1)\pi i\right)
=
-1.
\]
故每个 \(\omega_M^{-k_\ell}\) 都是 \(x^{F_d}+1\) 的根。并且
\[
\omega_M^{-k_\ell}
=
\exp\!\left(-\frac{(2\ell+1)\pi i}{F_d}\right)
\qquad (0\le \ell\le F_d-1)
\]
恰给出 \(x^{F_d}+1\) 的全部 \(F_d\) 个互异复根。因此 \(x^{F_d}+1\) 的全部根都是 \(D_m(x)\) 的根，从而
\[
x^{F_d}+1\mid D_m(x)
\qquad\text{于 }\CC[x].
\]
由于二者都属于 \(\ZZ[x]\)，Gauss 引理进一步给出
\[
x^{F_d}+1\mid D_m(x)
\qquad\text{于 }\ZZ[x].
\]

当 \(m+2\equiv 0\pmod 6\) 时，推论 \ref{cor:fold-zero-half-index-multiple6} 恰给出
\[
F_{(m+2)/2}\mid \frac{F_{m+2}}{2},
\]
故特殊情形立即成立。证毕。
\end{proof}

\noindent
上述链条把表 \ref{tab:fold_multiplicity_histogram} 中的隐藏位、容量缺口与碰撞指纹严格并入统一的信息几何结构；同时又把 \(\log C_m\) 的渐近模态提升为一条逐阶可反演、可任意阶消模并伴随非 \(C\)-有限障碍与尾部刚性的 Bernoulli--\(\zeta\) 谱链，并把约数驱动的谱零块最终刚化为群环多项式上的显式整系数因子。

\endinput
